\documentclass{article}

\usepackage{amsfonts}
\usepackage[T1]{fontenc}
\usepackage[french]{babel}
\begin{document}
\tableofcontents

\section{Définition}
Fonction Booléenne:
$$ f : \mathbb{F}_2^n \to \mathbb{F}_2 $$

Soit $f$ une fonction booléenne en $n$ variable on définie le poid de $f$ comme :
$$ \mathrm{wt}(f) = \#\{f(x) = 1 | x \in \mathbb{F}_2^n\} $$

Soit $f$ une fonction booléenne en $n$ variable, on peut la représenter par son vecteur d'évaluation (table de vérité) comme :
$$ (f(0), f(1), ..., f(2^n-1)) $$

On définie la distance de hamming de deux fonctions booléenne $f$ et $ g$ comme :
$$ d_H = \#\{f(x) \neq g(x) | x \in \mathbb{F}_2^n\}$$
ou
$$ \mathrm{wt}(f+g) $$

On définie le code de Reed-Muller en $n$ variable et de degrée $k$ $\mathrm{RM}(k, n)$ comme :
$$ \{f \in \mathbb{F}_2^n | deg(f) \leq k\} $$

Soit $C$ un Code $\subset \mathbb{F}_2^{2^n}$ et $f$ une fonction booléenne en $n$ variable on définie sa distance au code comme :
$$d(f, C) = \min_{c \in C} d_H(f, c)$$

\section{Runtime}
\subsection{Algo probabiliste}
\begin{tabular}{|c | c | c|}
	\hline
	nb itérations & nombre de fonctions & temp d'éxécution (seconde) \\
	\hline\hline

	10000 & 19424327 & 18633 \\
	100000 & 10107 & 188 \\
	100000 & stab 4: 11101773 & 14000 \\
	100000 & stab 6: 5354539 & 7000 \\
	100000 & stab 7: 201428 & 258 \\
	100000 & stab 8: 5290982 & 6670 \\
	\hline
\end{tabular}

\section{NL 2 de bent 8}
\begin{tabular}{|c|c|c|c|}
	\hline
	& 88 & 92 & 96 \\
	\hline\hline
	stab 1 & 0 & 0 & 0 \\
	\hline
	stab 2 & 0 & 0 & 0 \\
	\hline
	stab 4 & 0 & 0 & 0 \\
	\hline
	stab 6 & 0 & 0 & 0 \\
	\hline
	stab 7 & 0 & 0 & 0 \\
	\hline
	stab 8 & 0 & 0 & 0 \\
	\hline
\end{tabular}


\end{document}
